% !TeX program = xelatex
\documentclass[11pt,UTF8,fontset=fandol]{ctexart}

\usepackage[margin=0.82in]{geometry}
\usepackage{amsmath,amssymb}
\usepackage{booktabs,tabularx,array}
\usepackage{float}
\usepackage{needspace}
\usepackage[dvipsnames]{xcolor}
\usepackage{enumitem}
\usepackage{xurl}
\usepackage{hyperref}

\hypersetup{
  colorlinks=true,
  linkcolor=MidnightBlue,
  citecolor=ForestGreen,
  urlcolor=BrickRed,
  pdftitle={基于离线交互历史的个性化提示词优化}
}

\newcolumntype{Y}{>{\raggedright\arraybackslash}X}
\newcolumntype{L}[1]{>{\raggedright\arraybackslash}p{#1}}
\renewcommand{\arraystretch}{1.16}
\setlength{\tabcolsep}{5pt}
\setlist[itemize]{leftmargin=1.4em,itemsep=2pt,topsep=3pt}
\setlist[enumerate]{leftmargin=1.6em,itemsep=3pt,topsep=3pt}

\newcommand{\history}{D_{\mathrm{hist}}}
\newcommand{\dopt}{D_{\mathrm{opt}}}
\newcommand{\dselect}{D_{\mathrm{select}}}
\newcommand{\dtest}{D_{\mathrm{test}}}
\newcommand{\prompt}{p_u}
\newcommand{\agent}{\operatorname{FA}}
\newcommand{\judge}{\operatorname{Judge}}
\newcommand{\optimizer}{\operatorname{Revise}}

\begin{document}

\begin{center}
  {\LARGE\bfseries 基于离线交互历史的个性化提示词优化\par}
  \vspace{5pt}
  {\large 冻结宿主模型，只优化每位用户的可读执行提示词\par}
  \vspace{4pt}
  {\small 方法设计笔记 \quad 2026 年 8 月 29 日}
\end{center}

\section{结论}

可以写一个与“离线策略学习”对应的提示词版本。它不训练本地策略模型，也不修改宿主模型。它从同一位用户的历史中生成、测试并改写一份提示词：
\begin{equation}
  x
  \longrightarrow
  \agent(\prompt,x)
  \longrightarrow y.
\end{equation}

但不能声称“此前没有人做过提示词版本”。已有工作至少覆盖了三块：
\begin{itemize}
  \item DSO 从历史点击等反馈学习个性化提示词策略。策略根据用户和问题选择候选提示词，宿主模型保持冻结 \cite{kiyohara2025prompt}。
  \item AGP 自动改写用于推荐重排的用户画像生成提示词 \cite{wang2025agp}。
  \item GEPA 根据执行轨迹和文字反馈迭代提示词 \cite{agrawal2025gepa}。
\end{itemize}

因此，新的研究点不能只是“用提示词做个性化”或“自动优化提示词”。更具体的问题是：
\begin{quote}
能否从一个用户的多会话智能体历史中，优化一份可读、可编辑的执行提示词，使它在后续任务中更符合该用户，同时不降低任务成功率？
\end{quote}

下面给出一个可测试的方案。它还不是经过实验验证的算法。

\section{问题设定}

\subsection{目标}

给定用户历史 \(\history\)，学习一份文本提示词 \(\prompt\)。宿主智能体保持冻结。提示词只控制三类执行行为：
\begin{itemize}
  \item 何时使用当前会话的上下文；
  \item 何时直接回答、核实信息或先追问；
  \item 回答应直接、简短，还是分步展开。
\end{itemize}

提示词不替代任务推理。工具选择、代码执行、检索和最终回答仍由宿主智能体完成。

\Needspace{0.38\textheight}
\subsection{与离线策略版的区别}

\begin{table}[H]
  \centering
  \caption{两个版本优化的对象不同。}
  \begin{tabularx}{\textwidth}{@{}L{0.20\textwidth}L{0.35\textwidth}Y@{}}
    \toprule
    项目 & 离线策略版 & 提示词版 \\
    \midrule
    学习对象 &
    动作策略 \(\pi_\theta(a\mid x)\) &
    一份文本提示词 \(\prompt\) \\
    是否训练参数 &
    训练本地策略模型 &
    不训练宿主或本地策略模型 \\
    推理方式 &
    先选离散动作，再编译成指令 &
    宿主直接读取带条件的提示词 \\
    用户能否检查 &
    只能间接检查动作分布 &
    可以直接阅读、修改和删除 \\
    表达能力 &
    受固定动作空间限制 &
    可写条件规则，但可能变长或互相冲突 \\
    \bottomrule
  \end{tabularx}
\end{table}

\subsection{设计原则}

\begin{itemize}
  \item \textbf{宿主冻结：}所有候选提示词使用同一个宿主模型和同一套工具。
  \item \textbf{证据可查：}新增规则必须能指向训练侧的真实会话。
  \item \textbf{任务优先：}个性化提高时，任务成功率不能下降。
  \item \textbf{最终可读：}部署对象是一份普通文本，不是隐藏向量或软提示词。
  \item \textbf{测试隔离：}最终测试会话不参与提示词生成、改写或选择。
\end{itemize}

\section{角色与符号}

\begin{table}[H]
  \centering
  \caption{系统中的角色。}
  \begin{tabularx}{\textwidth}{@{}L{0.22\textwidth}L{0.15\textwidth}Y@{}}
    \toprule
    角色 & 是否更新 & 职责 \\
    \midrule
    宿主智能体 \(\agent\) &
    否 &
    读取提示词，完成任务并生成最终回答 \\
    提示词改写器 &
    否 &
    根据失败案例和文字反馈提出新提示词 \\
    评审模型 \(\judge\) &
    否 &
    比较回答是否符合训练侧用户证据，并给出简短修改建议 \\
    环境检查器 &
    否 &
    在可以自动验证的任务上判断是否真正完成任务 \\
    候选库 &
    更新文本 &
    保存提示词、来源、分数和每次改动 \\
    \bottomrule
  \end{tabularx}
\end{table}

\begin{table}[H]
  \centering
  \caption{主要符号。}
  \begin{tabularx}{0.94\textwidth}{@{}L{0.25\textwidth}Y@{}}
    \toprule
    符号 & 含义 \\
    \midrule
    \(\history\) & 同一用户的完整历史 \\
    \(\dopt\) & 生成初始提示词和改写反馈的数据 \\
    \(\dselect\) & 选择最终提示词的数据 \\
    \(\dtest\) & 只用于最终评测的数据 \\
    \(p^{(t)}\) & 第 \(t\) 轮使用的提示词 \\
    \(y(p,x)\) & 宿主在提示词 \(p\) 下对任务 \(x\) 的输出 \\
    \(\mathcal{P}_t\) & 第 \(t\) 轮结束后的候选提示词库 \\
    \bottomrule
  \end{tabularx}
\end{table}

\section{提示词的结构}

提示词沿用原策略版的三个执行维度，但把离散动作改成条件规则：
\begin{equation}
  p=(p_{\mathrm{context}},p_{\mathrm{information}},
     p_{\mathrm{response}},p_{\mathrm{avoid}}).
\end{equation}

\begin{table}[H]
  \centering
  \caption{一份提示词包含四段。}
  \begin{tabularx}{\textwidth}{@{}L{0.23\textwidth}Y@{}}
    \toprule
    段落 & 内容 \\
    \midrule
    上下文使用 &
    何时参考当前会话，何时只看本轮问题，以及哪些历史不能使用 \\
    信息获取 &
    何时直接回答，何时先核实，何时必须追问缺失信息 \\
    回答方式 &
    哪些任务需要直接结论、短答或分步说明 \\
    避免事项 &
    用户反复拒绝的行为，以及不能覆盖的宿主安全规则 \\
    \bottomrule
  \end{tabularx}
\end{table}

每条规则写成“条件 + 行为”。例如：
\begin{quote}
如果用户要求证明或推导，保留关键中间步骤；如果用户只要求结果，直接给出答案。
\end{quote}
不要只写“用户喜欢详细回答”。前一种写法说明何时生效，也减少不同偏好之间的冲突。

\section{数据边界}

\subsection{按会话和时间划分}

先按时间排列完整会话，再按任务类别分层。默认可以使用 \(60\%/20\%/20\%\) 的优化、选择和测试划分。比例必须在运行前固定。会话不可拆开：
\begin{equation}
  \dopt \prec \dselect \prec \dtest.
\end{equation}

\begin{table}[H]
  \centering
  \caption{每部分数据的用途。}
  \begin{tabularx}{\textwidth}{@{}L{0.20\textwidth}Y@{}}
    \toprule
    数据 & 允许的用途 \\
    \midrule
    \(\dopt\) &
    提取用户证据、生成初始提示词，并在内部轮换反馈任务和检查任务 \\
    \(\dselect\) &
    搜索结束后比较候选库，并选择部署版本 \\
    \(\dtest\) &
    最终一次评测；不能进入改写器或评审模型的上下文 \\
    \bottomrule
  \end{tabularx}
\end{table}

优化任务优先使用 \(\dopt\) 中的真实问题。模拟问题只用于补足稀少类别。模拟问题必须过滤与 \(\dselect\) 和 \(\dtest\) 的近重复内容。

\subsection{用户证据}

从 \(\dopt\) 中提取反复出现、跨任务且仍然有效的偏好。每条证据保留会话编号和时间。改写器提出新规则时，必须列出支持它的证据编号。找不到证据的规则不进入提示词。

\section{初始提示词}

强推理模型先从 \(\dopt\) 生成结构化初稿：
\begin{equation}
  p^{(0)}=\operatorname{SummarizePreferences}(\dopt).
\end{equation}

初稿只负责整理历史，不做搜索。它是必要基线。若后续优化不能稳定超过这份初稿，就没有证据说明自动改写有帮助。

\section{离线提示词优化}

\subsection{生成执行结果和反馈}

第 \(t\) 轮从 \(\dopt\) 选择一个按类别平衡的小批量 \(B_t\)。对每个任务 \(x\)，使用同一宿主、温度和工具设置运行当前提示词：
\begin{equation}
  y^{(t)}_x=\agent\!\left(p^{(t)},x\right).
\end{equation}

评审模型读取：
\begin{itemize}
  \item 与任务相关的训练侧用户证据；
  \item 当前任务 \(x\)；
  \item 宿主输出 \(y^{(t)}_x\)；
  \item 环境检查结果（若有）。
\end{itemize}

它返回三项内容：个性化判断、任务完成判断和一句具体修改建议。修改建议必须指出哪条提示词规则造成了问题，或缺少哪条有证据支持的规则。

\subsection{改写候选}

改写器根据当前提示词和失败反馈生成少量候选。每个候选只改一个段落，并给出文本差异。这样可以知道是哪项修改产生了效果：
\begin{equation}
  \{p'_1,\ldots,p'_K\}
  =\optimizer\!\left(p^{(t)},B_t,\text{feedback}\right).
\end{equation}

候选必须满足三个硬条件：
\begin{itemize}
  \item 不加入训练历史中没有支持的用户偏好；
  \item 不删除宿主原有的安全和任务规则；
  \item 不超过预先设定的提示词长度。
\end{itemize}

\subsection{选择规则}

每轮从 \(\dopt\) 另取检查批次 \(V_t\)，并保证 \(V_t\cap B_t=\varnothing\)。候选不能在产生其反馈的同一批任务上给自己打分。各轮轮换 \(B_t\) 和 \(V_t\)，分别计算个性化效果和任务成功率：
\begin{align}
  U_{\mathrm{pref}}(p;V_t)
    &=\frac{1}{|V_t|}
      \sum_{x\in V_t}
      J_{\mathrm{pref}}\!\left(\dopt,x,y(p,x)\right),\\
  U_{\mathrm{task}}(p;V_t)
    &=\frac{1}{|V_t|}
      \sum_{x\in V_t}
      J_{\mathrm{task}}\!\left(x,y(p,x)\right).
\end{align}

候选 \(p'\) 只有同时满足下面两项，才能替换当前提示词 \(p\)：
\begin{equation}
  U_{\mathrm{pref}}(p';V_t)>U_{\mathrm{pref}}(p;V_t),
  \qquad
  U_{\mathrm{task}}(p';V_t)\ge U_{\mathrm{task}}(p;V_t).
\end{equation}

还要按类别检查结果。不能用一个总体平均分掩盖某类任务的明显下降。若两个候选各自在不同类别更好，就都保留在候选库里，等更多选择数据后再决定。

\subsection{停止}

达到预先设定的运行预算，或连续若干轮没有候选通过选择规则时停止。搜索停止后，才首次用 \(\dselect\) 比较候选库并选择 \(p^\star\)。整个过程不读取 \(\dtest\)。

\section{完整流程}

\begin{enumerate}
  \item 按会话、时间和任务类别划分 \(\dopt,\dselect,\dtest\)。
  \item 从 \(\dopt\) 提取带证据编号的用户偏好。
  \item 生成结构化初始提示词 \(p^{(0)}\)。
  \item 在 \(\dopt\) 任务上运行宿主并收集失败反馈。
  \item 每轮只改一个提示词段落，生成候选。
  \item 在 \(\dopt\) 内部轮换检查批次，比较候选。
  \item 只接受个性化更好且任务成功率不下降的候选。
  \item 搜索结束后，用 \(\dselect\) 选择 \(p^\star\)。
  \item 最后在 \(\dtest\) 上评测一次。
\end{enumerate}

\section{推理}

部署时不再调用改写器或评审模型。系统只加载宿主和最终提示词：
\begin{equation}
  y=\agent\!\left(p^\star,x\right).
\end{equation}

用户可以查看、修改或删除 \(p^\star\)。代价是每次请求都要携带这份文本，因此必须单独报告提示词长度和输入 token 成本。

\section{评测}

\subsection{最少需要的对照}

\begin{table}[H]
  \centering
  \caption{建议的对照系统。}
  \begin{tabularx}{\textwidth}{@{}L{0.31\textwidth}Y@{}}
    \toprule
    系统 & 回答的问题 \\
    \midrule
    无个性化提示词 &
    宿主自身能做到什么？ \\
    直接提取的用户画像提示词 &
    只整理历史是否已经足够？ \\
    本文的优化提示词 &
    离线改写是否超过直接提取？ \\
    通用 GEPA 优化 &
    改进是否来自已有提示词优化器？ \\
    DSO 提示词策略（有日志概率时） &
    动态选择提示词是否优于一份固定提示词？ \\
    原离线策略版 &
    文本提示词与参数化动作策略哪一个更有效？ \\
    \bottomrule
  \end{tabularx}
\end{table}

\subsection{分开报告的指标}

\begin{itemize}
  \item 后续真实任务中的个性化符合度；
  \item 环境可验证的任务成功率；
  \item 各任务类别的结果和宏平均；
  \item 提示词长度、离线调用次数、推理 token 和延迟；
  \item 不同时间段上的偏好稳定性。
\end{itemize}

不要把这些指标压成一个总分。置信区间用于报告结果，不直接决定每次提示词改写。

\section{已有工作与创新边界}

\begin{table}[H]
  \centering
  \caption{最接近的工作已经覆盖了哪些部分。}
  \begin{tabularx}{\textwidth}{@{}L{0.20\textwidth}L{0.32\textwidth}Y@{}}
    \toprule
    工作 & 已经做到 & 与本方案的差别 \\
    \midrule
    OPPU \cite{tan2024oppu} &
    从每位用户的历史构造画像提示词，并与检索和参数训练比较 &
    不迭代搜索一份执行提示词 \\
    DSO \cite{kiyohara2025prompt} &
    从日志反馈学习按用户和问题选择候选提示词的策略 &
    优化的是提示词选择策略；主要实验是个性化电影描述，不是多步智能体任务 \\
    AGP \cite{wang2025agp} &
    自动优化用于推荐重排的用户画像生成提示词 &
    优化对象是推荐场景中的画像生成，不是通用智能体执行规则 \\
    GEPA \cite{agrawal2025gepa} &
    用执行轨迹和文字反馈进化提示词 &
    默认优化任务程序，不专门建模单个用户的长期偏好 \\
    本方案 &
    从一位用户的多会话历史优化一份可读执行提示词 &
    同时检查个性化和真实任务成功率 \\
    \bottomrule
  \end{tabularx}
\end{table}

因此，下面这些说法不成立：
\begin{itemize}
  \item “首次用提示词做个性化”；
  \item “首次从离线反馈优化提示词”；
  \item “首次在冻结模型上学习提示词”。
\end{itemize}

可能尚未被直接覆盖的是更窄的组合：\textbf{单个用户、多会话智能体历史、一份可读执行提示词、任务成功约束、未来真实任务评测}。目前只能把它称为待验证的研究差异。正式投稿前仍需做更完整的文献检索。

\section{建议先做的实验}

\begin{enumerate}
  \item 先使用现有用户历史，按会话和时间建立三段数据。
  \item 固定宿主模型、工具、温度和最大输出长度。
  \item 先跑“无提示词”和“直接提取提示词”。
  \item 再实现上述单段改写和选择循环。
  \item 用相同预算运行 GEPA，避免把已有优化器的能力算成新方法。
  \item 与原来的 18 动作策略版做同题比较。
  \item 只有在真实测试任务上同时改善个性化和任务成功，才继续做部署后的在线更新。
\end{enumerate}

{\footnotesize\raggedright
\begin{thebibliography}{9}
\setlength{\itemsep}{0pt}

\bibitem{kiyohara2025prompt}
Haruka Kiyohara et al.
\newblock Prompt Optimization with Logged Bandit Data.
\newblock \href{https://arxiv.org/abs/2504.02646}{arXiv:2504.02646}, 2025.

\bibitem{wang2025agp}
Chen Wang et al.
\newblock Automating Personalization: Prompt Optimization for Recommendation
Reranking.
\newblock \href{https://arxiv.org/abs/2504.03965}{arXiv:2504.03965}, 2025.

\bibitem{agrawal2025gepa}
Lakshya A. Agrawal et al.
\newblock GEPA: Reflective Prompt Evolution Can Outperform Reinforcement
Learning.
\newblock \href{https://arxiv.org/abs/2507.19457}{arXiv:2507.19457}, 2025.

\bibitem{tan2024oppu}
Zhaoxuan Tan et al.
\newblock Democratizing Large Language Models via Personalized
Parameter-Efficient Fine-Tuning.
\newblock \href{https://aclanthology.org/2024.emnlp-main.372/}{EMNLP}, 2024.

\end{thebibliography}
}

\end{document}
